% CONTENIDO MAESTRO: MÁQUINAS ELÉCTRICAS II
% Este archivo contiene la esencia técnica genérica para ser adaptada a cualquier institución.

\section*{FUNDAMENTACIÓN}
Esta propuesta formativa plantea la apropiación de las características constructivas y los fenómenos electromagnéticos que dan vida a las Máquinas Eléctricas y al funcionamiento de la maquinaria eléctrica más empleada en el sector técnico, ya sea como motores o generadores eléctricos. Se analizan los motores asincrónicos trifásicos, con sus principios de funcionamiento, conceptualizando comportamientos energéticos de rotores y representándolos a través de diagramas vectoriales y circuitos equivalentes. Como saberes eléctricos se destacan las curvas que hacen referencia a los arranques de los motores eléctricos, ya sea de forma directa o estrella-triángulo, dependiendo de la potencia eléctrica instalada. También se analizarán los tipos de arranque resistivo, reactivo, con empleo diverso de capacitores, etc., al observar el funcionamiento de las máquinas monofásicas en particular.

\section*{MATRIZ DE SABERES}

% EJE 1
\subsection*{EJE 1: FUNDAMENTOS DE MÁQUINAS MONOFÁSICAS}
\textbf{Saberes Prioritarios:}
\begin{itemize}
    \item Principios básicos de máquinas eléctricas AC.
    \item Leyes electromagnéticas y campos AC.
    \item Frecuencia y velocidad de sincronismo.
    \item Electrodinámica en máquinas industriales.
    \item Ángulos geométricos y eléctricos.
    \item FEM inducida, factores de arrollamiento y paso.
    \item Motores paso a paso.
\end{itemize}
\textbf{Estrategias:}
\begin{itemize}
    \item Modelos físicos y virtuales de motores asíncronos.
    \item Análisis de diagramas y videos de construcción.
    \item Herramientas visuales y experimentales para ángulos eléctricos.
    \item Ejemplos concretos de inducción electromagnética.
\end{itemize}

% EJE 2
\subsection*{EJE 2: MOTORES ASINCRÓNICOS TRIFÁSICOS}
\textbf{Saberes Prioritarios:}
\begin{itemize}
    \item Naturaleza eléctrica de máquinas asincrónicas.
    \item Aspectos constructivos y resbalamiento.
    \item Motor como transformador energético.
    \item Diagramas vectoriales y circuitos equivalentes.
    \item Curvas de arranque, cuplas y balance de potencias.
    \item Formas de arranque trifásico y frenado eléctrico.
\end{itemize}
\textbf{Estrategias:}
\begin{itemize}
    \item Simuladores de funcionamiento y partes del motor.
    \item Actividades de medición, ensamblaje y cálculo de resbalamiento.
    \item Prácticas de laboratorio sobre curvas de arranque.
    \item Uso de simulaciones para diagramas vectoriales.
\end{itemize}

% EJE 3
\subsection*{EJE 3: ENSAYOS EN MÁQUINAS SINCRÓNICAS}
\textbf{Saberes Prioritarios:}
\begin{itemize}
    \item Funcionamiento de generadores y motores sincrónicos.
    \item Máquinas en vacío y en carga.
    \item Reacción de inducido y regulación.
    \item Sincronización y puesta en paralelo industrial.
    \item Corrección del factor de potencia.
    \item Fuerza contraelectromotriz.
\end{itemize}
\textbf{Estrategias:}
\begin{itemize}
    \item Resolución de problemas y ejemplos prácticos.
    \item Videos y simulaciones de puesta en paralelo.
    \item Ensayos reales en Laboratorio Eléctrico.
    \item Recursos audiovisuales sobre fuerza contraelectromotriz.
\end{itemize}